\documentclass[10pt]{article}
\usepackage[margin=0.8in]{geometry}
\usepackage{amsmath,amssymb,amsthm,mathtools,booktabs,microtype,hyperref}
\hypersetup{hidelinks,pdftitle={Fixed square decompositions away from a singular contact}}
\newtheorem{proposition}{Proposition}
\newcommand{\R}{\mathbb R}
\title{Fixed square decompositions away from a singular contact}
\author{}\date{}
\begin{document}
\maketitle\vspace{-2.5em}

The coefficients of a Bell inequality may enter a Gram matrix, while its
testing functions must stay fixed. The cancellation below achieves this
separation on two complementary regions of the measurement space. It
combines the two state directions before completing squares, so every
testing square root has a fixed argument.

Write the source as
$\mathbf c=(d,\lambda,n,m,C,A,D,\alpha,\beta)\in\R_+^9$.
The compact testing coordinates satisfy
\[
 a^2+b^2=c^2+e^2=1,\qquad b\geq0.
\]
For finite angles they are $(a,b)=(u,1)/\sqrt{1+u^2}$ and
$(c,e)=(v,1)/\sqrt{1+v^2}$. The circle coordinates also include the
projective endpoints. All four real state components $W,X,Y,Z$ are free.
Define the state-linear expressions
\begin{align*}
 E&=aW-bY,& K&=aX-bZ,& J&=eE+cK,& U&=eZ-cY,\\
 E_\lambda&=eW+cX,& L_A&=bX+aZ,&
 M&=b(cW-eX)+a(cY-eZ).
\end{align*}
The compact canonical Bell slack is
\begin{equation}
 P=\lambda E_\lambda^2+\beta L_A^2+nX^2+mM^2
   +CE^2+AY^2+(D+d/2)J^2+(\alpha+d/2)U^2
   +\frac d2(L_A^2-Z^2).
 \label{eq:parent}
\end{equation}
Thus $P=F_h/4$, with the compact homogeneous normalization of the
canonical parent. The only negative term is explicit.

\section{A fixed family on an Alice arc}

Fix $0<B<1$ and restrict the tests to $0\leq b\leq B$.
For a source scalar $\tau$ put
\[
 S_\tau=\begin{pmatrix}C-\tau&\tau\\ \tau&D-\tau\end{pmatrix},
 \qquad \rho=\tau-d/2,\qquad
 h=\sqrt{\frac{2\alpha}{2\alpha+d}},\qquad
 \vartheta=\frac h{1+h}.
\]
Set $h=0$ if $\alpha=d=0$, and write $P_B=B-\vartheta B^2$.

\begin{proposition}
If
\begin{equation}
 S_\tau\succeq0,\qquad \rho\geq0,\qquad
 4A\rho\geq d^2P_B^2,                         \label{eq:admission}
\end{equation}
then $P$ has a positive semidefinite Gram representation on the following
twenty fixed testing functions:
\begin{align}
 \mathcal H_B=\big(&E_\lambda,L_A,X,M,cE,eE,J,R,ab\ell,U,b\ell,\ell,
 \nonumber\\[-2pt]
 &b\sqrt{1-b}\,\ell,\ b(1-b)\ell,\ cX,Y_0,bV,bY_0,
 \sqrt{b(B-b)}V,\sqrt{b(B-b)}Y_0\big),        \label{eq:features}
\end{align}
where $R=beX+aU$, $\ell=cY$, $V=cK$, and $Y_0=eY$.
Its matrix is affine in the nine source coordinates and three free
auxiliaries $\tau,p,q$.
\end{proposition}

\begin{proof}
Since $J-eE=cK$, the exact identity
\[
 CE^2+DJ^2=C(cE)^2+(eE,J)S_\tau(eE,J)^T+\tau(cK)^2
\]
keeps the large source coefficients in one positive block. Also
$L_A^2-Z^2=X^2-K^2$. Removing the first four event squares in
\eqref{eq:parent}, $C(cE)^2$, the $S_\tau$ block and $dJ^2/2$ leaves
\begin{equation}
 G=\frac d2R^2+\alpha U^2+A\ell^2+\frac d2b^2\ell^2
       +\frac d2(cX)^2+\rho V^2+AY_0^2+dbVY_0.
 \label{eq:residual}
\end{equation}
The coupling between $V$ and $Y_0$ can be redistributed because
\begin{equation}
 bVY_0=aR\ell-U\ell-b^2\ell^2.              \label{eq:cancellation}
\end{equation}
Using the fraction $\vartheta b$ of this identity gives
\begin{align}
G={}&\frac d2(R+\vartheta ab\ell)^2
 +\left(\sqrt\alpha\,U-
       \sqrt{\frac{d(1-h)}{2(1+h)}}\,b\ell\right)^2
 +A\ell^2+\frac d2(cX)^2 \nonumber\\
 &+\frac{dh}{(1+h)^2}(b\sqrt{1-b}\,\ell)^2
 +\frac{dh^2}{2(1+h)^2}(b(1-b)\ell)^2
 +\rho V^2+AY_0^2+d(b-\vartheta b^2)VY_0. \label{eq:squares}
\end{align}
The relation $2\alpha(1-h^2)=dh^2$ verifies this identity,
including $d=0$ and $\alpha=0$.

It remains to complete the last two-variable form. Its matrix has the
fixed interpolation
\begin{equation}
 \begin{pmatrix}\rho&d(b-\vartheta b^2)/2\\
 d(b-\vartheta b^2)/2&A\end{pmatrix}
 =\frac{b^2}{B^2}H_+
  +\frac{(B-b)^2}{B^2}H_0+b(B-b)Q_*,
 \label{eq:interpolation}
\end{equation}
where
\[
 H_+=\begin{pmatrix}\rho&dP_B/2\\dP_B/2&A\end{pmatrix},
 \qquad H_0=\begin{pmatrix}\rho&0\\0&A\end{pmatrix},
 \qquad Q_*=\begin{pmatrix}2\rho/B^2&d/(2B)\\d/(2B)&2A/B^2\end{pmatrix}.
\]
Condition \eqref{eq:admission} makes $H_+$ and $H_0$ positive
semidefinite. Since $0\leq\vartheta\leq1/2$, $P_B\geq B/2$, and hence
$\det Q_*=4A\rho/B^4-d^2/(4B^2)\geq0$ as well.
All factors in \eqref{eq:squares}--\eqref{eq:interpolation} belong to
the span \eqref{eq:features}. Their only testing radicals are
$\sqrt{1-b}$ and $\sqrt{b(B-b)}$.

To obtain an affine pencil, expand the displayed Gram and replace
$d\vartheta$ by $p$ and $d\vartheta^2$ by $q$.
For example the two square blocks on $(R,ab\ell)$ and $(U,b\ell)$ are
\[
 \begin{pmatrix}d/2&p/2\\p/2&q/2\end{pmatrix},
 \qquad
 \begin{pmatrix}\alpha&-p/2\\-p/2&d/2-p\end{pmatrix}.
\]
The other coefficients in \eqref{eq:squares} are $p-q$ and $q/2$.
In \eqref{eq:interpolation}, the off-diagonal of $H_+$ becomes
$(dB-pB^2)/2$. Each coefficient is now affine in
$(\mathbf c,\tau,p,q)$, and direct substitution of the testing functions
cancels all three auxiliary coefficients. Thus
\[
 P=\mathcal H_B^TQ_{20}(\mathbf c,\tau,p,q)\mathcal H_B
 \quad\hbox{for every }\tau,p,q\in\R.
\]
The nonlinear choices $p=d\vartheta$, $q=d\vartheta^2$ establish one
positive semidefinite witness; they impose no equations on the pencil.
\end{proof}

\section{The complementary Bob region}

Fix $\delta_0>0$ and restrict to $|c|\geq\delta_0$. Put
$\tau_0=d/(2\delta_0^2)$. The identity
\[
 CE^2+DJ^2-\frac d2K^2
 =C(cE)^2+(eE,J)S_{\tau_0}(eE,J)^T
    +\tau_0(c^2-\delta_0^2)K^2
\]
uses a third fixed radical. Substitution in \eqref{eq:parent} yields
\[
 P=\lambda E_\lambda^2+\beta L_A^2+(n+d/2)X^2+mM^2+AY^2
 +(\alpha+d/2)U^2+\frac d2J^2
 +C(cE)^2+(eE,J)S_{\tau_0}(eE,J)^T
 +\tau_0(c^2-\delta_0^2)K^2.
\]
Consequently $S_{\tau_0}\succeq0$ gives a source-affine Gram matrix on
\[
 \Psi=(E_\lambda,L_A,X,M,Y,U,J,cE,eE,
                     \sqrt{c^2-\delta_0^2}K).
\]
This matrix has order ten and no auxiliaries. Both constructions retain
all state signs and are valid at their testing boundaries. Adding the
nine scalar source-sign blocks gives total orders 29 and 19.

\section{A uniform native source domain}

The admission conditions can be verified without fixing the relative
rates of the source parameters. Normalize $d=1$ and
$\lambda=\beta=0$. Let
\[
 0<H<1,\quad 0<z,y\leq\epsilon H,\quad 0<r<1,\quad
 L=yt>0,\quad 0<t<a_s:=1+zr,\quad j=a_s-t.
\]
Write $B_s=H+L$, $T=B_s+zr$, $U_s=B_s-zyr$, $S=a_s+L$, and
\begin{align*}
 \Delta={}&z+(1+z)B_s+(1+z)B_s^2+2z(1+z)B_sr
       +z^3r^2+2z^2r\\
 &+y\{(1+z)B_s+z+z^2r(2-r)\},\\
 G_s={}&B_s+z(H+r-1),\\
 M_s={}&(1+z)TH+yG_s+yzj+zrL(1+y).
\end{align*}
The native C2 source coefficients are
\begin{align}
 n&=\frac{Ta_sS-r\Delta}{2r\Delta},&
 m&=\frac{z^2(1+y)(1-r)a_sS}{2G_s\Delta},&
 C&=\frac{a_sM_s}{2L\Delta},\nonumber\\
 A&=\frac{ya_sTS}{2H\Delta},&
 D&=\frac{(1+z)TU_sS}{2yj\Delta},&
 \alpha&=\frac{TM_s}{2(1-H)\Delta}.             \label{eq:native}
\end{align}
Require $n\geq0$. For $0<\epsilon\leq1/100$, the other displayed
denominators and coefficients are positive.

Here are bounds that remain uniform as any positive opening tends to
zero. Since $L<ya_s$, setting $q_\epsilon=1+\epsilon+\epsilon^2$ gives
\[
 H\leq B_s\leq q_\epsilon H,\quad
 1\leq a_s\leq1+\epsilon,\quad
 H\leq T\leq H(1+\epsilon)^2,\quad
 1\leq S\leq(1+\epsilon)^2,\quad U_s\geq H(1-\epsilon^2).
\]
In the positive polynomial for $\Delta$, substitution of these bounds
gives
\[
 0\leq\Delta-H(1+H)
 \leq\epsilon H(3+7H+2H^2)
 \leq6\epsilon H(1+H).
\]
For $M_s$, use $G_s\leq H(1+\epsilon)^2$ and $j\leq1+\epsilon$:
\[
 H^2\leq M_s
 \leq H^2(1+4\epsilon+7\epsilon^2+5\epsilon^3+\epsilon^4)
 \leq H^2(1+5\epsilon).
\]
These estimates first give $A\leq\epsilon$.
For $\eta=CD/(C+D)$ the reciprocal cancels the separate small openings:
\begin{align*}
 \frac1\eta
 &=2y\Delta\left\{\frac{t}{a_sM_s}
             +\frac{j}{(1+z)TU_sS}\right\}
 \leq\frac{4\epsilon(1+6\epsilon)}{1-\epsilon}
 \leq5\epsilon,\\
 \frac1{4A\eta}
 &=H\Delta^2\left\{\frac{t}{a_s^2TSM_s}
             +\frac{j}{(1+z)a_sT^2U_sS^2}\right\}
 \leq(1+H)^2\frac{(1+6\epsilon)^2}{1-\epsilon}.
\end{align*}
As $(1+6\epsilon)^2/(1-\epsilon)\leq1+14\epsilon$, it follows that
\[
 \eta\geq\frac1{5\epsilon},\qquad
 4A\eta\geq(1+H)^{-2}-14\epsilon.
\]
The square root entering the Gram witness is equally well controlled.
Put $\nu=TM_s/H^3$ and $\omega=\Delta/[H(1+H)]$.
Then $1\leq\nu\leq1+8\epsilon$, $1\leq\omega\leq1+6\epsilon$, and
\[
 \frac{h^2}{H^2}=
 \frac{\nu}{H^2\nu+(1-H^2)\omega}.
\]
Thus $H(1-3\epsilon)\leq h\leq H(1+4\epsilon)$, so
$|h-H|\leq4\epsilon$.

Now choose, once and for all,
\[
 \delta_0=2^{-26},\qquad B=1-\delta_0,\qquad
 \epsilon=\frac{\delta_0^2}{4096}=2^{-64}.
\]
Set $\tau=\eta-1/4$ and $\rho=\eta-3/4$.
For positive $C,D$, $S_\tau\succeq0$ whenever $0\leq\tau\leq\eta$;
indeed its determinant is $CD-\tau(C+D)$.
The bound on $\eta$ ensures this condition, $\rho>0$, and also
$S_{\tau_0}\succeq0$ for $\tau_0=1/(2\delta_0^2)$.

Finally, $h/(1+h)$ is $1$-Lipschitz on $[0,1]$. Replacing $h$ by
$H$ changes $P_B^2$ by at most $8\epsilon$. Therefore
\begin{align*}
 4A\rho-P_B^2
 &\geq(1+H)^{-2}
       -\left(B-\frac{H}{1+H}B^2\right)^2-25\epsilon\\
 &\geq\frac{\delta_0^2}{4}-25\epsilon
   >\frac{\delta_0^2}{8}.
\end{align*}
The second inequality uses
$B(1+H\delta_0)\leq1-\delta_0^2$ and $(1+H)^2\leq4$.
This proves both regional certificates on the entire displayed source
envelope. Arbitrary nonnegative $\lambda,\beta$ may then be restored as
their whole event squares. Literal $n=0$ is included. The rational
parametrization \eqref{eq:native} concerns positive openings; undefined
zero denominators are not assigned limiting values.


\section{Why the contact requires a different completion}

At an exact contact, every square in a positive decomposition must vanish.
For the native C2 contact parametrization, let
$z=u_c^2>0$, $y=v_c^2>0$, $0<H<1$ and $L=yt>0$. Its contact is
\[
 (u_c,v_c)=(\sqrt z,\sqrt y),\qquad
 (W_c,X_c,Y_c,Z_c)=
 \left(\frac{H+L}{\sqrt{zy}},-\sqrt z\,r,\frac H{\sqrt y},1\right).
\]
The native source equations give $P=0$ there, whereas
\[
 cE=\frac{L}{\sqrt{(1+z)(1+y)}}>0.
\]
For the source envelope above the contact lies in the omitted cap, since
\[
 b_c=(1+z)^{-1/2}\geq1-z/2>B,\qquad
 |c_c|=\sqrt{\frac{y}{1+y}}<\frac{\delta_0}{64}.
\]
For $C>0$, both displayed decompositions retain the strictly positive
term $C(cE)^2$. They therefore cannot remain positive semidefinite
through that contact. A completion covering it must redistribute this
term with the remaining state directions.

In general, if $P(t)=\Phi(t)^TQ\Phi(t)$ with $Q\succeq0$ and
$P(t_c)=0$, then $Q\Phi(t_c)=0$. This follows from
$P(t_c)=\|Q^{1/2}\Phi(t_c)\|^2$ and is a necessary constraint on any
replacement. It keeps the exact contact in the construction before
positivity is proved.
\end{document}
